% Based on: diffusion-for-failures paper.

\chapter{Diffusion Under Failure-Induced Embodiments}
\label{ch:failures}

% TODO: Apply the conditioning machinery from Chapter~\ref{ch:payload} to
% the actuation-failure problem first introduced in
% Section~\ref{sec:constrained:failures}.

\section{Revisiting the Long-Deployment Failure Problem}
\label{sec:failures:revisit}

% TODO: Recall the long-deployment, actuation-failure setting from
% Chapter~\ref{ch:constrained}, and the workspace consequences of locked
% joints.

\section{Conditioning on High-Level Encodings of Task Constraints}
\label{sec:failures:conditioning}

% TODO: Encode failed embodiments via joint-range constraints, exposed to
% the diffusion model as high-level conditioning.

\section{A Single Model Spanning Prehensile and Non-Prehensile Primitives}
\label{sec:failures:unified}

% TODO: One generative model handles both prehensile and non-prehensile
% primitives across the failure-induced embodiment space.

\section{Comparison with the Sampling-Based Approach}
\label{sec:failures:comparison}

% TODO: Compare against the sampling-based approach from
% Section~\ref{sec:constrained:failures} along the four design axes from
% Chapter~\ref{ch:framework}.
